\documentclass[11pt,a4paper]{article}

\usepackage{lmodern}
\usepackage[T1]{fontenc}
\usepackage[utf8]{inputenc}
\usepackage[margin=2.6cm]{geometry}
\usepackage{amsmath}
\usepackage{amssymb}
\usepackage{booktabs}
\usepackage{microtype}
\usepackage{xcolor}
\usepackage[round]{natbib}
\usepackage[hidelinks]{hyperref}

\newcommand{\code}[1]{\texttt{#1}}
\newcommand{\SR}{\mathrm{SR}}

\title{\code{jsharpe}: Sharpe-Ratio Inference in Practice, and Where the
Corrections Do Not Reach}

\author{Thomas Schmelzer\thanks{\url{https://github.com/tschm/jsharpe}}}

\date{\today}

\begin{document}

\maketitle

\begin{abstract}
\code{jsharpe} implements the standard toolkit for deciding whether a Sharpe
ratio is real: its asymptotic variance under non-Gaussian and autocorrelated
returns, the Probabilistic Sharpe Ratio, the minimum track record length,
critical values and power, and the multiple-testing machinery --- Bonferroni,
{\v S}id{\'a}k, Holm and a false-discovery-rate calibration --- that a
strategy-screening exercise needs. This paper documents what the package
computes and checks it against simulation and against the published formulas
it implements. Three things hold well and three do not. The variance estimator
is accurate to $3.4\%$ at $T = 100$ and $1.4\%$ at $T = 1000$; the FDR
calibration is internally exact, reproducing its own target $q$ to five
decimals; the FWER adjustments and the effective-rank statistic behave as
specified. Against that: the package's headline non-Gaussian correction is
\emph{exactly inert} at its own default null $\SR_0 = 0$, because the variance
is evaluated under the null, and we measure the cost --- a nominally $5\%$
test rejects $8.2\%$ of the time on severely skewed returns. The minimum track
record length omits both the leading $1$ and the observed-moment terms of the
published formula, understating the requirement by a factor of $4.6$ at
$\SR = 1$ with skew $-1.5$ and kurtosis $8$. And the trials parameter $K$
inverts the multiple-testing correction it is named for: raising it from $1$ to
$1000$ lowers the critical Sharpe ratio from $0.164$ to $0.058$, because two
functions in the package model the variance of a maximum of $K$ draws in
opposite directions --- one correctly (it falls with $K$), one as a
$1 + \log K$ inflation, off by a factor of $39$ at $K = 200$. Separately, the
cluster-count selector returns $k = n-1$ --- one pair and $58$ singletons --- on
a matrix with ten true blocks, because a singleton silhouette is scored $1$
rather than the conventional $0$. We give the one-line repairs for each and the
simulation harness that would have caught them.
\end{abstract}

\section{Introduction}

A Sharpe ratio \citep{sharpe1966,sharpe1994} is a ratio of two estimates, and
the literature on how badly that fact is usually ignored is now larger than
the literature on the statistic itself. \citet{lo2002} gave the asymptotic
distribution of the estimator, including under serial correlation;
\citet{mertens2002} and \citet{opdyke2007} extended it to non-Gaussian
marginals; \citet{bailey2012} recast the resulting test as a probability, the
Probabilistic Sharpe Ratio, and derived the track record length a claim
requires; \citet{bailey2014} and \citet{harvey2015} added the correction for
having looked at many strategies before picking one.

\code{jsharpe} is a small implementation of that toolkit --- roughly $1{,}500$
lines over six modules, depending only on NumPy \citep{harris2020} and SciPy
\citep{virtanen2020}. It is the kind of package whose value depends entirely on
correctness: the functions are short, their inputs are a handful of scalars,
and nobody inspects the output for plausibility because the output is the
plausibility check. That makes it worth testing against simulation rather than
against itself.

This paper does that. Section~\ref{sec:method} states the estimators as the
package computes them, in the notation of the sources, so the divergences later
are visible rather than asserted. Section~\ref{sec:results} tests each against
Monte Carlo or against the closed form it claims to implement. The results
split cleanly: the core variance estimator and the FDR calibration are right,
and three of the corrections that motivate the package do not reach the code
paths a user actually calls. Section~\ref{sec:fixes} states each repair as a
diff-sized change, because all three are.

\section{The estimators}
\label{sec:method}

Let $r_1, \ldots, r_T$ be returns with mean $\mu$, standard deviation $\sigma$,
skewness $\gamma_3$, non-excess kurtosis $\gamma_4$ and first-order
autocorrelation $\rho$, and let $\widehat{\SR} = \hat\mu / \hat\sigma$.

\subsection{Variance of the estimator}

The package computes
\begin{equation}
  V(\SR, T) = \frac{a - b \, \gamma_3 \SR + c \, \frac{\gamma_4 - 1}{4} \SR^2}{T}
  \cdot \mathrm{Var}[M_K],
  \label{eq:var}
\end{equation}
where $a = 1 + 2B$, $b = 1 + B + C$, $c = 1 + 2C$ with
$B = \rho/(1-\rho)$ and $C = \rho^2/(1-\rho^2)$ carrying the
autocorrelation adjustment, and $\mathrm{Var}[M_K]$ is the variance of the
maximum of $K$ standard normals, computed by Gauss--Hermite quadrature. At
$\rho = 0$ and $K = 1$, \eqref{eq:var} reduces to
$(1 - \gamma_3 \SR + \frac{\gamma_4-1}{4}\SR^2)/T$, which is the
\citet{bailey2012} form, and under Gaussian returns
($\gamma_3 = 0$, $\gamma_4 = 3$) to $(1 + \SR^2/2)/T$, the \citet{lo2002}
result. The $\mathrm{Var}[M_K]$ factor has no counterpart in those sources; it
is examined in Section~\ref{sec:maxk}.

\subsection{Tests built on it}

Every downstream quantity is a normal tail probability with
\eqref{eq:var} as its variance, evaluated \emph{under the null}:
\begin{align}
  \SR_c &= \SR_0 + z_{1-\alpha}\sqrt{V(\SR_0, T)}, \label{eq:crit}\\
  \mathrm{PSR} &= \Phi\!\left( \frac{\SR - \SR_0}{\sqrt{V(\SR_0, T)}} \right),
    \label{eq:psr}\\
  \mathrm{minTRL} &= V(\SR_0, 1) \left( \frac{z_{1-\alpha}}{\SR - \SR_0} \right)^2 ,
    \label{eq:mtrl}
\end{align}
with power $1 - \Phi((\SR_c - \SR_1)/\sqrt{V(\SR_1, T)})$. The choice to
evaluate the variance at $\SR_0$ rather than at the observed $\SR$ is the
hinge of Section~\ref{sec:inert}: it is defensible for \eqref{eq:crit} and
\eqref{eq:psr}, which are null-hypothesis statements, and it is a departure
from the published form of \eqref{eq:mtrl}, whose reference version is
\begin{equation}
  \mathrm{minTRL}^{\star} = 1 +
  \left( 1 - \gamma_3 \SR + \frac{\gamma_4 - 1}{4}\SR^2 \right)
  \left( \frac{z_{1-\alpha}}{\SR - \SR_0} \right)^2 .
  \label{eq:mtrlstar}
\end{equation}

\subsection{Screening many strategies}

For a pool of $M$ candidates the package offers the three standard
family-wise adjustments --- Bonferroni $\min(1, Mp)$, {\v S}id{\'a}k
$1 - (1-p)^M$ \citep{sidak1967}, and Holm's step-down procedure
\citep{holm1979} in either variant --- and a false-discovery-rate calibration
\citep{benjamini1995,storey2002} in the two-normal mixture form. In the latter,
a strategy is drawn from $\mathcal{N}(\SR_1, \sigma_1^2)$ with prior
probability $p_{H_1}$ and from $\mathcal{N}(\SR_0, \sigma_0^2)$ otherwise, and
the critical value solves
\begin{equation}
  \mathbb{P}[H_0 \mid X > c] =
  \left( 1 + \frac{\overline{\Phi}((c - \SR_1)/\sigma_1)}
                  {\overline{\Phi}((c - \SR_0)/\sigma_0)}
          \frac{p_{H_1}}{1 - p_{H_1}} \right)^{-1} = q ,
  \label{eq:fdr}
\end{equation}
with $\sigma_0^2 = V(\SR_0, T)$ and $\sigma_1^2 = V(\SR_1, T)$ and
$\overline{\Phi}$ the normal survival function. Separately, the expected best
of $K$ trials is the extreme-value approximation of \citet{bailey2014},
\begin{equation}
  \mathbb{E}[\max\nolimits_K] \approx \sigma \left[
    (1 - \gamma_e)\Phi^{-1}\!\left(1 - \tfrac{1}{K}\right)
    + \gamma_e \Phi^{-1}\!\left(1 - \tfrac{1}{Ke}\right) \right],
  \label{eq:emax}
\end{equation}
$\gamma_e$ the Euler--Mascheroni constant.

\subsection{Two utilities}

The package also carries a correlation-structure module and a portfolio
helper. \code{effective\_rank} is $\exp(H)$ for $H$ the entropy of the
normalised eigenvalues of a correlation matrix \citep{roy2007}, and
\code{number\_of\_clusters} selects $k$ by maximising the ratio of the mean
silhouette score \citep{rousseeuw1987} to its standard deviation, the
criterion of \citet{lopezdeprado2018}. \code{minimum\_variance\_weights\_%
for\_correlated\_assets} solves for the minimum-variance portfolio under the
assumption that the correlation matrix is \emph{equicorrelated} at the average
of its off-diagonal entries, so that the inverse is available in closed form by
Sherman--Morrison --- the constant-correlation structure of \citet{elton1976}.

\section{Implementation}
\label{sec:impl}

The package is pure functions over scalars and small arrays, with no state and
no configuration. Three details of that surface matter for what follows.

\paragraph{Every moment argument is optional, and defaults to Gaussian.}
$\gamma_3 = 0$, $\gamma_4 = 3$, $\rho = 0$, $K = 1$. A user who supplies none
of them gets the textbook Gaussian iid test, which is the right default. A user
who supplies them expects them to matter.

\paragraph{$\SR_0$ defaults to $0$ everywhere it appears.} It is the natural
null --- ``does this strategy make money'' --- and it is what the package's own
documentation and worked example use.

\paragraph{The quadrature is shared.}
$\mathrm{Var}[M_K]$ and $\mathbb{E}[M_K]$ come from a $200$-node
Gauss--Hermite rule constructed once at import. We checked it against
simulation: at $K = 10$ the quadrature gives
$\mathrm{Var}[M_{10}] = 0.3443$ against $0.3448$ from $4 \times 10^5$ draws,
and at $K = 200$, $0.1607$ against $0.1600$. The quadrature is not the
problem anywhere below.

\section{Results}
\label{sec:results}

\subsection{The variance estimator is sound}

\begin{table}[t]
\centering
\caption{Variance of $\widehat{\SR}$ from \eqref{eq:var} against $2 \times
10^4$ simulations of iid Gaussian returns with true $\SR = 0.5$.}
\label{tab:var}
\begin{tabular}{lrrr}
\toprule
$T$ & Simulated & Equation~\eqref{eq:var} & Error \\
\midrule
$24$   & $0.052147$ & $0.046875$ & $-10.1\%$ \\
$100$  & $0.011650$ & $0.011250$ & $-3.4\%$ \\
$1000$ & $0.001110$ & $0.001125$ & $+1.4\%$ \\
\bottomrule
\end{tabular}
\end{table}

Table~\ref{tab:var} confirms \eqref{eq:var} as an asymptotic result behaving
asymptotically. At $T = 1000$ it is accurate to $1.4\%$, at $T = 100$ to
$3.4\%$, and at $T = 24$ it understates the true variance by $10.1\%$. The
last figure deserves a note because $T = 24$ is the default in the package's
FDR routine and the order of magnitude of its documented examples: at two
years of monthly data, the asymptotic variance is a tenth too small and every
test built on it is correspondingly liberal, before any of the issues below.

\subsection{The non-Gaussian correction is inert under the default null}
\label{sec:inert}

The package's first advertised feature is the correction for skewness and
excess kurtosis. Substituting $\SR = \SR_0$ into \eqref{eq:var}, as
\eqref{eq:crit}--\eqref{eq:mtrl} all do, gives
\begin{equation}
  V(\SR_0, T)\big|_{\SR_0 = 0} = \frac{a}{T} \cdot \mathrm{Var}[M_K],
\end{equation}
in which $\gamma_3$ and $\gamma_4$ do not appear. At the default null they are
multiplied by zero. This is not an approximation that degrades; it is exact
cancellation.

\begin{table}[t]
\centering
\caption{The moment arguments at $\SR_0 = 0$, for $\SR = 0.5$ and $T = 100$.
All four rows are bit-identical. The last two columns repeat the calculation
at $\SR_0 = 0.3$, where the arguments do reach the result.}
\label{tab:inert}
\begin{tabular}{rrrrrrr}
\toprule
$\gamma_3$ & $\gamma_4$ & $V(\SR_0, T)$ & PSR & minTRL & PSR ($\SR_0{=}0.3$) & minTRL ($\SR_0{=}0.3$) \\
\midrule
$0.0$  & $3.0$  & $0.01000000$ & $0.9999997133$ & $10.822174$ & $0.974795$ & $70.68$ \\
$-1.5$ & $8.0$  & $0.01000000$ & $0.9999997133$ & $10.822174$ & $0.942654$ & $108.73$ \\
$0.0$  & $20.0$ & $0.01000000$ & $0.9999997133$ & $10.822174$ & --- & --- \\
$2.0$  & $3.0$  & $0.01000000$ & $0.9999997133$ & $10.822174$ & --- & --- \\
\bottomrule
\end{tabular}
\end{table}

Table~\ref{tab:inert} shows the cancellation at work. A user testing a
$\SR = 0.5$ track record of $T = 100$ observations against $\SR_0 = 0$ gets
the same PSR to ten decimal places whether the returns are Gaussian or have
kurtosis $20$. Shifting the null to $\SR_0 = 0.3$ restores the correction:
the PSR falls from $0.975$ to $0.943$ and the required track record grows by
$54\%$.

The autocorrelation argument is the exception. It enters through $a = 1 + 2B$,
which does not vanish at $\SR_0 = 0$, so $\rho$ works as documented at every
null: raising it from $0$ to $0.5$ lifts the minTRL at $\SR = 0.5$ from
$10.8$ to $32.5$ and the critical Sharpe ratio at $T = 100$ from $0.164$ to
$0.285$.

\paragraph{What the inertness costs.}
Whether it matters depends on whether the correction was doing anything. It
was.

\begin{table}[t]
\centering
\caption{Rejection rate of a nominally $5\%$ one-sided test of $H_0: \SR = 0$
when $H_0$ is true, over $2 \times 10^5$ series of $T = 100$ from the
package's own non-Gaussian generator. ``Package'' is \eqref{eq:psr} with the
observed moments supplied; ``observed-moment form'' evaluates the same moment
polynomial at the observed $\SR$ instead, and with $T-1$ in place of $T$, as
\citet{bailey2012} prescribe. Skewness and kurtosis are the sample averages.}
\label{tab:calib}
\begin{tabular}{lrrrr}
\toprule
Distribution & Mean skew & Mean kurt.\ & Package & Observed-moment form \\
\midrule
\code{gaussian} & $-0.001$ & $2.940$  & $5.04\%$ & $4.84\%$ \\
\code{mild}     & $-0.773$ & $5.088$  & $5.71\%$ & $4.90\%$ \\
\code{moderate} & $-1.410$ & $8.377$  & $7.26\%$ & $5.75\%$ \\
\code{severe}   & $-1.789$ & $11.567$ & $8.18\%$ & $6.86\%$ \\
\bottomrule
\end{tabular}
\end{table}

Table~\ref{tab:calib} is the practical statement of the problem. Under
Gaussian returns the test is calibrated: it rejects $5.04\%$ of the time at a
nominal $5\%$. As the returns acquire the left skew and fat tails that hedge
fund returns are notorious for, the test drifts liberal, reaching $8.18\%$ ---
$64\%$ more false discoveries than advertised --- on the generator's severe
setting. Evaluating the same moment polynomial at the observed Sharpe ratio,
which is the one-line change of Section~\ref{sec:fixes}, recovers about half
of the gap, to $6.86\%$. It does not close it, because at $T = 100$ the
sample skewness and kurtosis are themselves badly estimated, and no
plug-in correction can fix that.

\paragraph{The documented example.}
The package's own worked example is a $T = 36$ series with
$\widehat{\SR} = 0.8150$, sample skewness $-0.581$ and kurtosis $2.188$. Its
PSR is reported as effectively $1$, and it is: $0.999999$, identically with and
without the moment arguments. The published form gives $0.999904$ --- the same
conclusion. The minTRL is where the example diverges: the package returns
$4.07$ observations, \eqref{eq:mtrlstar} returns $7.81$, and the honest reading
of a $36$-month record is that the second number is the one to quote.

\subsection{The minimum track record length is not the published one}

\begin{table}[t]
\centering
\caption{\code{minimum\_track\_record\_length} against
\eqref{eq:mtrlstar}, $\SR_0 = 0$, $\alpha = 0.05$. The package's value does
not depend on the moments; the reference value does.}
\label{tab:mtrl}
\begin{tabular}{rrrrr}
\toprule
$\SR$ & $\gamma_3$ & $\gamma_4$ & Package & Equation~\eqref{eq:mtrlstar} \\
\midrule
$0.3$ & $0.0$  & $3.0$ & $30.06$ & $32.41$ \\
$0.3$ & $-1.5$ & $8.0$ & $30.06$ & $49.32$ \\
$0.5$ & $0.0$  & $3.0$ & $10.82$ & $13.17$ \\
$0.5$ & $-1.5$ & $8.0$ & $10.82$ & $24.67$ \\
$1.0$ & $0.0$  & $3.0$ & $2.71$  & $5.06$  \\
$1.0$ & $-1.5$ & $8.0$ & $2.71$  & $12.50$ \\
\bottomrule
\end{tabular}
\end{table}

Table~\ref{tab:mtrl} separates two divergences from \eqref{eq:mtrlstar}. The
first is the missing leading $1$, which costs one observation and matters only
where the answer is small --- but the answer \emph{is} small exactly where the
function is most likely to be quoted, and at $\SR = 1$ the package's $2.71$
against a reference $5.06$ is a factor of $1.9$ from that term alone. The
second is the inertness of Section~\ref{sec:inert}, which under a Gaussian
null removes the moment polynomial entirely: at $\SR = 1$ with skew $-1.5$ and
kurtosis $8$, the reference requirement is $12.50$ observations and the
package reports $2.71$, understating it by a factor of $4.6$.

A minimum track record length is a statement made to somebody else --- an
allocator, a risk committee --- about how much evidence exists. Understating
it by a factor of four in exactly the fat-tailed, negatively skewed case that
motivates computing it at all is the most consequential of the divergences in
this paper, and the cheapest to repair.

\subsection{The two models of a maximum of \texorpdfstring{$K$}{K} disagree}
\label{sec:maxk}

The package contains two functions concerned with the maximum of $K$ Sharpe
ratios. \code{moments\_Mk} computes $\mathbb{E}[M_K]$ and
$\mathrm{Var}[M_K]$ for $K$ standard normals by quadrature, and
\code{variance\_of\_the\_maximum\_of\_k\_Sharpe\_ratios} returns
$\sigma^2 (1 + \log K)$, documented as a ``logarithmic variance inflation
factor'' that ``increases monotonically with $K$'' because ``selection across
a larger pool increases the uncertainty of the selected estimate''.

\begin{table}[t]
\centering
\caption{Variance and expectation of the maximum of $K$ standard normals:
$4 \times 10^5$ simulations against the two package functions. The
quadrature is right; the $1 + \log K$ inflation has the wrong sign of
dependence on $K$.}
\label{tab:maxk}
\begin{tabular}{rrrrrrr}
\toprule
& \multicolumn{3}{c}{$\mathrm{Var}[M_K]$} & \multicolumn{3}{c}{$\mathbb{E}[M_K]$} \\
\cmidrule(lr){2-4}\cmidrule(lr){5-7}
$K$ & Simulated & \code{moments\_Mk} & $1 + \log K$ & Simulated & \code{moments\_Mk} & Eq.~\eqref{eq:emax} \\
\midrule
$1$   & $1.0028$ & $1.0000$ & $1.0000$ & $0.0002$ & $0.0000$ & --- \\
$2$   & $0.6815$ & $0.6817$ & $1.6931$ & $0.5649$ & $0.5642$ & $0.5198$ \\
$10$  & $0.3448$ & $0.3443$ & $3.3026$ & $1.5383$ & $1.5388$ & $1.5746$ \\
$50$  & $0.2159$ & $0.2157$ & $4.9120$ & $2.2489$ & $2.2491$ & $2.2763$ \\
$200$ & $0.1600$ & $0.1607$ & $6.2983$ & $2.7471$ & $2.7460$ & $2.7655$ \\
\bottomrule
\end{tabular}
\end{table}

Table~\ref{tab:maxk} settles it. The variance of the maximum of $K$ iid
normals \emph{falls} with $K$ --- the maximum of many draws is pinned near the
upper tail and is therefore less variable, a standard order-statistic result
\citep{david2003} --- from $1.00$ at $K = 1$ to $0.16$ at $K = 200$. The
quadrature reproduces this to three decimals. The $1 + \log K$ function
returns $6.30$ where the truth is $0.16$: wrong by a factor of $39$, and
wrong in direction. The prose in its docstring describes a real phenomenon ---
selection does make the selected estimate less trustworthy --- but that
phenomenon lives in $\mathbb{E}[M_K]$, not $\mathrm{Var}[M_K]$, and
\eqref{eq:emax} already captures it, agreeing with simulation to within
$0.02$ at $K \geq 10$.

\paragraph{Which makes $K$ a discount, not a penalty.}
The consequence falls on the tests, because \eqref{eq:var} --- not the
$1+\log K$ function --- is what $K$ actually reaches, and it enters as a
\emph{multiplier} $\mathrm{Var}[M_K] < 1$.

\begin{table}[t]
\centering
\caption{Effect of the trials parameter $K$ on the tests, at $T = 100$,
$\SR_0 = 0$, $\alpha = 0.05$. Larger $K$ makes every test easier to pass.}
\label{tab:kdirection}
\begin{tabular}{rrrr}
\toprule
$K$ & Critical $\SR$ & Power at $\SR_1 = 0.3$ & PSR at $\SR = 0.2$ \\
\midrule
$1$    & $0.164485$ & $0.9075$ & $0.977250$ \\
$10$   & $0.096521$ & $0.9881$ & $0.999673$ \\
$100$  & $0.070634$ & $0.9990$ & $0.999998$ \\
$1000$ & $0.057794$ & $0.9999$ & $1.000000$ \\
\bottomrule
\end{tabular}
\end{table}

Table~\ref{tab:kdirection} is the trap. A user who screened a thousand
strategies, read the documentation for $K$ --- ``number of strategies for
multiple testing adjustment'' --- and passed $K = 1000$ gets a critical Sharpe
ratio of $0.058$ instead of $0.164$, and a PSR of $1.000000$ for an observed
$0.2$. The parameter named for the multiple-testing problem makes the
multiple-testing problem worse. It is not that the correction is too weak; its
sign is reversed.

The fix is not to change $\mathrm{Var}[M_K]$, which is correct as a variance.
It is that a deflation for selection shifts the null --- replace $\SR_0$ by
$\SR_0 + \mathbb{E}[\max_K]$, which is what \eqref{eq:emax} exists for and
what \citet{bailey2014} prescribe --- rather than rescaling the variance in
place. Used that way, the two ingredients the package already has combine
correctly; used as $K$, they cancel the correction and then reverse it.

\subsection{What holds}

\begin{table}[t]
\centering
\caption{FDR calibration at $\SR_0 = 0$, $\SR_1 = 0.5$, $p_{H_1} = 0.05$,
$T = 24$. \emph{$\alpha$ direct} recomputes
$\overline{\Phi}((\SR_c - \SR_0)/\sigma_0)$ from the returned threshold;
$\hat q$ is the routine's own re-derivation of the target.}
\label{tab:fdr}
\begin{tabular}{rrrrrr}
\toprule
$q$ & $\SR_c$ & $\alpha$ & $\alpha$ direct & Power & $\hat q$ \\
\midrule
$0.05$ & $0.6563$ & $0.00065$ & $0.00065$ & $0.2351$ & $0.05000$ \\
$0.10$ & $0.5872$ & $0.00201$ & $0.00201$ & $0.3435$ & $0.10000$ \\
$0.25$ & $0.4792$ & $0.00944$ & $0.00944$ & $0.5382$ & $0.25000$ \\
$0.50$ & $0.3597$ & $0.03903$ & $0.03903$ & $0.7415$ & $0.50000$ \\
$0.75$ & $0.2181$ & $0.14267$ & $0.14267$ & $0.9036$ & $0.75000$ \\
\bottomrule
\end{tabular}
\end{table}

Table~\ref{tab:fdr} checks the FDR routine two ways: the significance level it
reports is recovered exactly from the threshold it returns, and its own
re-derivation of $q$ matches the requested $q$ to five decimals across the
range. The mixture in \eqref{eq:fdr} is solved by bisection on
$[-10, 10]$ with the two degenerate cases --- no root in the interval, and no
root at all when $\sigma_0 \gg \sigma_1$ with small $q$ --- returning
$-\infty$ and NaN respectively rather than a spurious number. The trade-off it
exposes is the real one: controlling the false-discovery rate at $5\%$ against
a $5\%$ prior leaves the test with only $24\%$ power.

The family-wise adjustments are textbook and behave so. On
$p = (0.001, 0.008, 0.02, 0.04, 0.30)$: Bonferroni gives
$(0.005, 0.040, 0.100, 0.200, 1.000)$, {\v S}id{\'a}k the uniformly smaller
$(0.005, 0.039, 0.096, 0.185, 0.832)$, and Holm the smaller still
$(0.005, 0.032, 0.060, 0.080, 0.300)$, monotone in the sorted order and never
below the raw values. The effective rank recovers block structure as it
should: on a $60$-strategy correlation matrix built from $10$ latent factors
with noise $0.1$ it returns $6.81$, and it rises smoothly with the noise
($6.51$ at noise $0.05$, $7.83$ at $0.2$, $13.9$ at $0.5$, $28.7$ at $1.0$,
$49.7$ at $3.0$) towards the full dimension, rather than jumping. It
under-counts the latent factors --- the blocks are unequal in size and the
weakest contribute little entropy --- but it is monotone in the right
direction and needs no partition to compute.

The autocorrelated generator delivers roughly what it promises. Requesting
$\rho = 0.2$ yields a measured $0.171$ and $\rho = 0.5$ yields $0.440$, the
shortfall being the expected consequence of applying a nonlinear marginal
transform to a Gaussian AR(1): the copula's rank dependence survives, the
Pearson autocorrelation is compressed by $12$ to $15\%$. The target marginals
themselves come through --- skewness around $-2.4$ and kurtosis around $16$
for the severe setting, at all three $\rho$.

\subsection{The constant-correlation portfolio, quantified}

\begin{table}[t]
\centering
\caption{\code{minimum\_variance\_weights\_for\_correlated\_assets} against
the exact minimum-variance solution, on one-factor correlations
$C_{ij} = b_i b_j$ with loadings drawn uniformly from the stated range and
volatilities from $[0.08, 0.30]$. $2000$ draws per row; ``penalty'' is the
mean excess of the achieved variance over the minimum.}
\label{tab:mv}
\begin{tabular}{rlrrrr}
\toprule
$n$ & Loadings & Mean $\rho$ & SD$(\rho)$ & Mean $\max_i |\Delta w_i|$ & Variance penalty \\
\midrule
$3$  & $[0.55, 0.65]$ & $0.360$ & $0.013$ & $0.0076$ & $0.03\%$ \\
$3$  & $[0.30, 0.90]$ & $0.360$ & $0.079$ & $0.0544$ & $1.43\%$ \\
$5$  & $[0.55, 0.65]$ & $0.361$ & $0.018$ & $0.0152$ & $0.14\%$ \\
$5$  & $[0.30, 0.90]$ & $0.360$ & $0.111$ & $0.1018$ & $5.91\%$ \\
$10$ & $[0.55, 0.65]$ & $0.360$ & $0.021$ & $0.0185$ & $0.43\%$ \\
$10$ & $[0.30, 0.90]$ & $0.358$ & $0.132$ & $0.1281$ & $17.80\%$ \\
$25$ & $[0.55, 0.65]$ & $0.360$ & $0.023$ & $0.0139$ & $0.99\%$ \\
$25$ & $[0.30, 0.90]$ & $0.360$ & $0.144$ & $0.1005$ & $38.67\%$ \\
$25$ & $[0.00, 0.95]$ & $0.226$ & $0.190$ & $0.1171$ & $79.77\%$ \\
\bottomrule
\end{tabular}
\end{table}

The equicorrelation assumption is documented, so Table~\ref{tab:mv} is a
calibration rather than a criticism. It is exact where it claims to be: for
$n = 2$, where the average off-diagonal correlation \emph{is} the only
correlation, and for genuinely equicorrelated matrices at any $n$, both to
machine precision. Away from that, the cost is governed by the dispersion of
the correlations and grows with $n$: at $n = 25$ with correlations spread
around $0.36$ with a standard deviation of $0.14$, the portfolio it returns
carries $39\%$ more variance than the minimum. It is a shrinkage estimator
without the shrinkage-intensity parameter --- \citet{ledoit2004} with
$\lambda$ pinned at $1$ --- and given how sensitive exact mean--variance
solutions are to their inputs \citep{best1991,markowitz1952}, that is a
defensible default and a bad surprise. We note it because the docstring says
``assumes'' where the results say ``costs $39\%$''. Neither this function nor
the exact solution is long-only: both hold short positions in essentially
every draw at $n \geq 10$.

\subsection{The cluster-count criterion degenerates}

\begin{table}[t]
\centering
\caption{\code{number\_of\_clusters} on correlation matrices generated with a
known block structure ($T = 200$ observations, noise $0.1$), matrix seed $0$
and $k$-means seed $1$. The last two columns are the quality score at the true
$k$ and at the near-singleton partition $k = n - 1$.}
\label{tab:clusters}
\begin{tabular}{rrrrrr}
\toprule
$n$ & True $k$ & Effective rank & Selected $k$ & Quality at true $k$ & Quality at $k = n-1$ \\
\midrule
$20$  & $4$  & $2.93$ & $3$  & $1.20$ & $6.61$  \\
$40$  & $5$  & $4.41$ & $39$ & $2.06$ & $69.00$ \\
$60$  & $10$ & $6.81$ & $59$ & $1.42$ & $68.51$ \\
$100$ & $10$ & $9.06$ & $99$ & $1.45$ & $7.31$  \\
\bottomrule
\end{tabular}
\end{table}

Table~\ref{tab:clusters} shows the criterion failing in a specific and
diagnosable way. For $n = 60$ with ten true blocks it returns $k = 59$: a
partition of one pair and $58$ singletons. The quality profile explains it ---
$3.87$ at $k = 2$, $3.34$ at $k = 5$, $1.42$ at the true $k = 10$, a trough of
$0.71$ at $k = 20$, then a climb through $4.26$ at $k = 58$ to $68.51$ at
$k = 59$. The selection is not an artefact of one $k$-means initialisation:
$k = 59$ is chosen for all five initialisation seeds we tried. There are two
problems compounding.

The first is the criterion. The mean silhouette divided by its standard
deviation is unbounded: it diverges as the spread of the silhouette scores
goes to zero, whatever the mean is, so a partition in which all points score
alike wins regardless of whether the scores are good. The implementation
guards against exactly zero spread and nothing more.

The second makes the first reachable. In the silhouette computation, a
singleton cluster is given $a(i) = 0$, so its score is
$s(i) = (b(i) - 0)/b(i) = 1$ --- the maximum possible. The convention since
\citet{rousseeuw1987} is $s(i) = 0$ for a singleton, and it exists precisely
to stop a partition into singletons from scoring perfectly. With $58$ of $60$
points scoring exactly $1$, the spread collapses and the ratio explodes. The
effective rank on the same matrices, which needs no partition at all, returns
$6.81$ and $9.06$ against a true $10$.

Only the smallest case escapes, and not by much: at $n = 20$ the criterion
picks $k = 3$ against a true $4$, with the degenerate partition the runner-up
at $6.61$ against $7.02$. From $n = 40$ upwards the degenerate partition wins
outright. The magnitudes in the last column vary with the $k$-means seed --- a
quality of $69$ is a near-zero denominator, and near-zero denominators are
noisy --- but the selected $k$ does not.

\section{Repairs}
\label{sec:fixes}

All three substantive findings are small changes.

\begin{enumerate}
  \item \textbf{Evaluate the moment polynomial at the observed Sharpe ratio}
        in \eqref{eq:psr} and \eqref{eq:mtrl}, keeping the null variance for
        \eqref{eq:crit} where it belongs. This is a change of argument, from
        \code{sharpe\_ratio\_variance(SR0, ...)} to
        \code{sharpe\_ratio\_variance(SR, ...)}, and it moves the calibration
        of Table~\ref{tab:calib} from $8.18\%$ to $6.86\%$.
  \item \textbf{Add the leading $1$} to \eqref{eq:mtrl}, making it
        \eqref{eq:mtrlstar}. One character, a factor of $1.9$ at $\SR = 1$.
  \item \textbf{Do not route $K$ through the variance.} Either delete the
        $\mathrm{Var}[M_K]$ factor from \eqref{eq:var} and apply selection as
        a shift of $\SR_0$ by \eqref{eq:emax}, or document $K$ as what it
        currently is, which is not a multiple-testing correction. Replace the
        $1 + \log K$ function with \code{moments\_Mk(K)[2]}, which the package
        already computes correctly.
  \item \textbf{Score singletons at zero} in the silhouette, and bound the
        cluster-quality criterion --- or cap \code{max\_clusters} well below
        $n - 1$. The first is the standard convention; the second is the
        cheaper guard.
\end{enumerate}

None of these is hard to test. The simulations in Tables~\ref{tab:var},
\ref{tab:calib} and \ref{tab:maxk} took seconds to run, and each of them
fails loudly against the current code. A package whose entire output is a
calibration claim is a package whose test suite should contain a calibration
check: draw from a known null, count the rejections, assert the count is near
nominal. That test would have caught findings one and three, and its absence
is why a documented feature could be exactly inert without anything going red.

\section{Limitations}

Every number here is a simulation or a closed form, so the usual caveat about
sample specificity does not apply --- but the choice of null distributions
does. The calibration in Table~\ref{tab:calib} uses the package's own
two-component Gaussian mixture, which reaches skewness $-1.8$ and kurtosis
$11.6$ in sample at $T = 100$. Real hedge fund returns are often more extreme
and, more importantly, serially dependent in ways an iid mixture is not; we did
not test calibration under autocorrelation, where the $\rho$ adjustment is
active and may behave differently.

The comparison of \eqref{eq:mtrl} against \eqref{eq:mtrlstar} treats the latter
as the reference. That is a choice of authority, not a proof: \eqref{eq:mtrlstar}
is the form in \citet{bailey2012}, and we have checked its internal consistency
with \eqref{eq:psr} rather than its own calibration.

Table~\ref{tab:mv} uses one-factor correlations, which are a well-conditioned
and optimistic family. Empirical correlation matrices estimated from short
samples are neither, and the variance penalties there would be larger and
noisier. The clustering results use one generator with equal-sized latent
factors; the degeneracy of Section~\ref{sec:results} is a property of the
criterion rather than of that generator, but the $n$ at which it takes over is
not.

Finally, this paper says nothing about the two functions in the package that
are conditional on inputs we did not vary systematically --- the observed-FDR
routine and the power calculation --- beyond noting that both are algebraic
consequences of \eqref{eq:psr} and \eqref{eq:crit} and therefore inherit
Section~\ref{sec:inert} exactly.

\section{Conclusion}

The half of \code{jsharpe} that computes distributions is right: the
asymptotic variance matches simulation to a few percent for $T \geq 100$, the
order-statistic quadrature is accurate to three decimals, the extreme-value
approximation for the expected best of $K$ trials tracks simulation within
$0.02$, the FDR calibration is internally exact, and the FWER adjustments and
effective rank are textbook.

The half that turns those distributions into a verdict has three defects, and
they share a shape: a correction is available, is documented, is passed by the
user, and does not reach the answer. Skewness and kurtosis are exactly
cancelled by evaluating the variance under a zero null, and the cost is a
nominally $5\%$ test rejecting $8.2\%$ of the time on fat-tailed returns. The
minimum track record length drops the same terms and the leading constant,
understating a fat-tailed requirement fourfold. And the trials parameter, which
exists to punish having searched, rewards it instead, because a variance of a
maximum was used where a shift of a mean was needed.

The pattern is worth naming, because it is not specific to this package. Each
defect is invisible to a unit test that checks a function returns a float in
$[0, 1]$, and each is caught immediately by a test that draws from a known null
and counts rejections. For a library whose whole purpose is to say how often we
should expect to be wrong, that is the test that matters.

\appendix

\section{Reproducibility}
\label{app:repro}

Everything above runs against the package in this repository with NumPy and
SciPy. The inertness of Section~\ref{sec:inert}:

\begin{verbatim}
from jsharpe import probabilistic_sharpe_ratio, minimum_track_record_length
for g3, g4 in [(0.0, 3.0), (-1.5, 8.0), (0.0, 20.0), (2.0, 3.0)]:
    print(g3, g4,
          probabilistic_sharpe_ratio(0.5, 0.0, T=100, gamma3=g3, gamma4=g4),
          minimum_track_record_length(0.5, 0.0, gamma3=g3, gamma4=g4))
\end{verbatim}

Table~\ref{tab:calib}, with \code{np.random.seed(0)} and the package's own
generator:

\begin{verbatim}
import numpy as np, scipy.stats as st
from jsharpe import probabilistic_sharpe_ratio, generate_non_gaussian_data

def bldp(SR, T, g3, g4):
    return st.norm.cdf(SR * np.sqrt(T - 1) / np.sqrt(1 - g3*SR + (g4-1)/4*SR**2))

for name in ("gaussian", "mild", "moderate", "severe"):
    X = generate_non_gaussian_data(200_000, 100, SR0=0.0, name=name)
    sr = X.mean(axis=1) / X.std(axis=1, ddof=1)
    g3, g4 = st.skew(X, axis=1), st.kurtosis(X, axis=1, fisher=False)
    pkg = [probabilistic_sharpe_ratio(float(a), 0.0, T=100, gamma3=float(b),
                                      gamma4=float(c))
           for a, b, c in zip(sr[:20000], g3[:20000], g4[:20000])]
    print(name, np.mean(np.array(pkg) > 0.95),
          np.mean(bldp(sr[:20000], 100, g3[:20000], g4[:20000]) > 0.95))
\end{verbatim}

Table~\ref{tab:maxk}:

\begin{verbatim}
from jsharpe.sharpe.quadrature import moments_Mk
from jsharpe import variance_of_the_maximum_of_k_Sharpe_ratios as vmax
rng = np.random.default_rng(0)
for K in (1, 2, 10, 50, 200):
    m = rng.standard_normal((400_000, K)).max(axis=1)
    print(K, m.var(ddof=1), moments_Mk(K)[2], vmax(K, 1.0), m.mean(), moments_Mk(K)[0])
\end{verbatim}

Table~\ref{tab:mv} draws $C_{ij} = b_i b_j$ with
\code{b = rng.uniform(lo, hi, n)} and unit diagonal, scales by
$\sigma \sim U[0.08, 0.30]$, and compares against
\code{np.linalg.solve(V, ones)} normalised to sum to one.
Table~\ref{tab:clusters} uses
\code{get\_random\_correlation\_matrix(n, k)} under
\code{np.random.seed(0)} followed by
\code{number\_of\_clusters(C, retries=10)}, whose second return value is the
quality dictionary the table's last two columns read.

\bibliographystyle{plainnat}
\bibliography{references}

\end{document}
